\subsection{Generalization to an Unseen Participant}
\label{subsec:unseen-participant}

We evaluated the models without adaptation on an unseen participant comprising
50 open-hand and 50 closed-hand trials. All 100 trials were accepted by the
preprocessing pipeline. Metrics were pooled over valid frames rather than
averaged across the two conditions, thereby accounting for their slightly
different durations. The recurrent baselines contained no gripper classifier:
IMU models received only inertial measurements, whereas fusion models used
early concatenation of sEMG and IMU features. GRU and LSTM widths were selected
to approximate the trainable parameter count of the proposed Stage-II model.

\begin{table*}[t]
 \centering
 \caption{Motion generalization on the unseen participant. Values are mean
 Euclidean errors pooled across open and closed trials. The proposed model uses
 its dense trajectory head at 100--200 ms and its intent head at 500--1000 ms.}
 \label{tab:unseen-motion}
 \small
 \setlength{\tabcolsep}{5pt}
 \begin{tabular}{lrrrrrr}
  \hline
  Model & Current (cm) & Pixel (px) & 100 ms (cm) & 200 ms (cm) &
  500 ms (cm) & 1000 ms (cm) \\
  \hline
  Proposed       & \textbf{7.29} & \textbf{409.7} & \textbf{7.63} &
                   \textbf{8.04} & 10.78 & 15.16 \\
  IMU--GRU       & 7.43 & 429.0 & 7.82 & 8.33 & \textbf{10.73} & 14.88 \\
  sEMG+IMU--GRU  & 11.40 & 498.7 & 11.38 & 11.73 & 13.34 & 16.14 \\
  IMU--LSTM      & 7.92 & 416.7 & 8.34 & 8.78 & 11.09 & \textbf{14.40} \\
  sEMG+IMU--LSTM & 8.88 & 448.8 & 9.16 & 9.83 & 11.79 & 15.18 \\
  \hline
 \end{tabular}
\end{table*}

The proposed model produced the lowest current-position and screen errors and
the lowest errors at 100 and 200 ms. Relative to the strongest conventional
baseline at each of these outputs, it reduced current-position error by
0.15 cm, pixel error by 7.0 pixels, 100-ms error by 0.19 cm, and 200-ms error
by 0.28 cm. At longer horizons, the advantage did not persist: IMU--GRU was
0.05 cm better at 500 ms and IMU--LSTM was 0.77 cm better at 1 s. Thus, the
unseen-participant evidence supports improved current and short-horizon
decoding, but does not establish superiority over recurrent IMU baselines at
all long horizons.

Naive early fusion generalized poorly, particularly for the GRU. This result
suggests that inter-participant variation in sEMG can corrupt a recurrent
motion representation when the modalities are concatenated without control.
In contrast, the proposed model restricts sEMG to a bounded correction of an
IMU-derived mechanical representation. Its stronger short-horizon performance
is therefore consistent with the intended role of sEMG as incremental evidence
rather than an unconstrained replacement for inertial information.

\begin{table}[t]
 \centering
 \caption{Proposed-model future error and causal reference baselines on the
 combined unseen open and closed trials.}
 \label{tab:unseen-future-baselines}
 \small
 \begin{tabular}{rrrr}
  \hline
  Horizon & Proposed & Persistence & Linear extrapolation \\
  \hline
  100 ms  & \textbf{7.63} & 8.57 & 8.19 \\
  200 ms  & \textbf{8.04} & 10.72 & 10.13 \\
  500 ms  & \textbf{10.78} & 19.34 & 21.81 \\
  1000 ms & \textbf{15.16} & 37.57 & 50.96 \\
  \hline
 \end{tabular}
\end{table}

The persistence baseline repeats the model's current Cartesian estimate,
whereas linear extrapolation estimates velocity by least-squares fitting over
the preceding 200 ms and propagates that velocity causally. At 1 s, the
proposed intent head reduced error by 22.41 cm relative to persistence and by
35.80 cm relative to linear extrapolation. This demonstrates predictive motion
information beyond holding the current pose or extending recent velocity,
although the matched IMU--LSTM remained more accurate at that horizon.

The proposed gripper classifier was assessed only after combining both
conditions, because classwise F1 is not meaningful within an all-open or
all-closed subset. The combined confusion matrix was
\begin{equation}
 \begin{bmatrix}
  10330 & 130\\
  759 & 10280
 \end{bmatrix},
\end{equation}
with rows denoting true open and closed states. This corresponds to 95.86\%
accuracy and 95.86\% macro-F1, indicating balanced transfer of the frozen
state representation to the unseen participant.
